\documentclass[tikz,border=8pt]{standalone}
\usepackage{tikz}
\usepackage{amsmath}
\usepackage{pgfplots}
\pgfplotsset{compat=1.18}
\usepackage{ctex}
\usetikzlibrary{calc, positioning}

\begin{document}
\begin{tikzpicture}

\begin{axis}[
  name=ax,
  width=11cm, height=7cm,
  xmin=-4.5, xmax=4.5,
  ymin=-0.02, ymax=0.50,
  xlabel={$x$},
  ylabel={$f(x)$},
  axis lines=left,
  xtick={-4,-3,-2,-1,0,1,2,3,4},
  ytick={0,0.1,0.2,0.3,0.4},
  yticklabel style={font=\scriptsize},
  xticklabel style={font=\scriptsize},
  grid=major, grid style={gray!30},
  clip=false,
  legend style={at={(1.02,1.0)}, anchor=north west, font=\small,
                fill=white, draw=gray!50, row sep=2pt}
]

  % N(0,1) 黑色实线
  \addplot[black, ultra thick, domain=-4.5:4.5, samples=200]
    {exp(-x^2/2)/sqrt(2*3.14159)};
  \addlegendentry{$N(0,1)$}

  % t(2) 蓝色虚线（重尾，峰低）
  % t(df) PDF: Γ((df+1)/2) / (sqrt(df*pi)*Γ(df/2)) * (1+x^2/df)^(-(df+1)/2)
  % df=2: 常数 = 1/(2*sqrt(2)) ≈ 0.3536，但直接用数值近似
  \addplot[blue!70!black, thick, dashed, domain=-4.5:4.5, samples=200]
    {0.3536*(1+x^2/2)^(-1.5)};
  \addlegendentry{$t(2)$}

  % t(5) 绿色点虚线
  % df=5: 常数 ≈ 0.3796
  \addplot[green!60!black, thick, densely dashed, domain=-4.5:4.5, samples=200]
    {0.3796*(1+x^2/5)^(-3.0)};
  \addlegendentry{$t(5)$}

  % t(30) 紫色点线
  % df=30: 常数 ≈ 0.3983（接近正态）
  \addplot[violet, thick, dotted, domain=-4.5:4.5, samples=200]
    {0.3983*(1+x^2/30)^(-15.5)};
  \addlegendentry{$t(30)$}

  % 峰值标注
  \node[font=\scriptsize, black] at (axis cs:0.5, 0.42) {$N(0,1)$峰值};
  \draw[->, thin, black] (axis cs:0.5, 0.41) -- (axis cs:0.1, {exp(0)/sqrt(2*3.14159)+0.01});

  % 说明注释
  \node[font=\scriptsize, blue!70!black, right] at (axis cs:1.8, 0.28) {自由度↑ $\Rightarrow$ 尾部↓};

\end{axis}

% 公式框（axis 外）
\node[font=\small, align=center, fill=white, draw=gray!50, thin,
      rounded corners=2pt, inner sep=5pt]
  at ($(ax.south)+(0,-1.1cm)$)
  {$t$ 分布：$f(x;\nu)=\dfrac{\Gamma\!\left(\tfrac{\nu+1}{2}\right)}%
    {\sqrt{\nu\pi}\,\Gamma\!\left(\tfrac{\nu}{2}\right)}%
    \left(1+\dfrac{x^2}{\nu}\right)^{-\frac{\nu+1}{2}}$
    \quad 当 $\nu\to\infty$ 时趋近 $N(0,1)$};

% 整体标题
\node[font=\large\bfseries] at ($(ax.north)+(0,0.6cm)$)
  {$z$ 分布 vs $t$ 分布（不同自由度）};

\end{tikzpicture}
\end{document}
